\section{Language Scope}
\label{sec:background}

This section introduces the key characteristics of the languages to which type slicing can be applied.

\subsection{Bidirectional Types}
\label{sec:bidirectional-types}

Type slicing is formalised for \textit{bidirectional type systems}~\cite{BidirectionalTypes}. Such a type system is an algorithmic specification of typing judgements, naturally lending itself to direct typing algorithms which do not require \textit{guessing} of types during type checking. These systems do require annotations, but annotation burden is lessened by locally inferring types~\cite{LocalInference}.

This is achieved by specifying the \textit{mode} of the type parameter in a typing judgement, distinguishing when it is an \textit{input} (type analysis/checking) and when it is an \textit{output} (type synthesis):
\[\synthesis{e}{\tau}\]
Read as: \textit{$e$ synthesises type $\tau$ under typing assumptions $\Gamma$ and type variable assumptions $\Delta$}. The type $\tau$ is an \textit{output}.
\[\analysis{e}{\tau}\]
Read as: \textit{$e$ analyses against type $\tau$ under typing assumptions $\Delta, \Gamma$}. The type $\tau$ is an \textit{input}. This is where the terminology of \textit{synthesis / analysis} slices is derived from.

Such languages have two fundamental rules:
\[\inference[\textsc{SVar}]{x : \tau \in \Gamma}{\synthesis{x}{\tau}} \qquad
  \inference[\textsc{Sub}]{\synthesis{e}{\tau}}{\analysis{e}{\tau}}\]
Variables synthesise their type from the typing assumptions. Subsumption bridges the two modes: any term that can synthesise a type also analyses against that same type.

\subsection{Downwards Static Graduality}
\label{sec:gradual-types}

To calculate a type slice we will need a way to \textit{gradually} omit sub-expressions from the program, and type check these less \textit{`precise'} partial programs at less or equally `precise' types. Formally, precision must form two partial orders (notated $\sqsubseteq$) over types and expressions, which can be lifted onto typing assumptions $\Gamma$.

\begin{definition} A bidirectional calculus satisfies \textit{downwards static graduality} if, for all $\Delta, \Gamma', \Gamma, e', e, \tau$ such that $\Gamma' \sqsubseteq \Gamma$ and $e' \sqsubseteq e$, both:
\begin{itemize}
  \item If $\synthesis{e}{\tau}$ then there exists $\tau' \sqsubseteq \tau$ such that $\synthesis[\Delta;\ \Gamma']{e'}{\tau'}$.
\item For all $\tau' \sqsubseteq \tau$, if $\analysis{e}{\tau}$ then $\analysis[\Delta; \Gamma']{e'}{\tau'}$.
\end{itemize}
\end{definition}

For example, this precision order can correspond to the precision relation on types in a gradually typed system \cite{GradualFunctional}. In this situation this graduality theorem corresponds to the downwards static case\footnote{Hence the name.} of the graduality theorem in the refined criteria for gradual types \cite{GradualRefined}, adapted to a bidirectional setting. The calculus we demonstrate type slicing on in this paper goes even further by allowing maximally incomplete programs, omitting any program sub-expression, based upon the Hazel calculus~\cite{HazelLivePaper}.

Although not necessary, it is also useful that precision forms a lattice with meets and joins to perform optimisations to type slicing calculation (\S\ref{sec:algorithms}).

%%% Local Variables:
%%% mode: LaTeX
%%% TeX-master: "main"
%%% End:
